\chapter{Thực nghiệm và đánh giá}

Chương này được tổ chức theo \emph{khẳng định} chứ không theo thứ tự các đợt chạy. Bảng
\ref{bang:lotrinh} nối mỗi khẳng định với mục chứng minh nó, để người đọc theo dõi được lập
luận mà không cần nhớ lịch sử thực nghiệm.

\begin{table}[htbp]
\centering
\caption{Lộ trình chương: mỗi khẳng định và nơi nó được kiểm chứng}
\label{bang:lotrinh}
\small
\begin{tabular}{@{}p{9.6cm}l@{}}
\toprule
\textbf{Khẳng định} & \textbf{Mục}\\
\midrule
Phần lớn xáo trộn mà hành khách chịu là do hệ thống tự gây ra, không phải do đường sá. & 4.2\\
Trần cam kết cưỡng chế được, và cái giá của nó đo được theo từng mức chặt chẽ. & 4.3\\
Cơ sở hành vi luôn khả thi; cơ sở kết quả có thể vô nghiệm và không bao giờ hơn. & 4.4, 4.7\\
Kết luận không phụ thuộc lựa chọn thước đo. & 4.5\\
Hình thức của lời bảo đảm chỉ có nghĩa nếu bộ sinh phương án thực thi được nó. & 4.6\\
Thước đo của văn liệu mù ở chế độ cam kết. & 4.8\\
Giá của lời hứa tăng theo mức tự do, trên toàn thang ngân sách. & 4.9\\
Kết quả không phải sản phẩm phụ của việc cắt ngắn tìm kiếm. & 4.11\\
\bottomrule
\end{tabular}
\end{table}

\section{Thiết lập thực nghiệm}
\label{muc:thietlap}

\subsection{Dữ liệu và bộ dữ liệu thử}

Dữ liệu là bộ chuyến đi taxi công khai của Manhattan do Taxi \& Limousine Commission của New
York phát hành, đã được chuẩn hóa và đóng gói trong bộ FleetPy \cite{fleetpy2022} theo giấy
phép CC-BY. Từ bộ này, khóa luận dựng một panel gồm 20 ô thực nghiệm: năm ngày làm việc từ 14
đến 18 tháng 11 năm 2018, mỗi ngày bốn lần lấy mẫu độc lập.

\begin{table}[htbp]
\centering
\caption{Tham số của mỗi ô thực nghiệm}
\label{bang:thamso}
\small
\begin{tabular}{@{}lr@{}}
\toprule
\textbf{Tham số} & \textbf{Giá trị}\\
\midrule
Khung thời gian & 08:00--10:00 giờ New York\\
Số yêu cầu & 108\\
Số xe & 8\\
Sức chứa mỗi xe & 4 chỗ\\
Cửa sổ đón $\Omega$ & 600 giây\\
Ràng buộc thời gian trên xe & $1{,}5 \times$ thời gian đi thẳng\\
Số nút mạng đường & 96\\
Hạt giống & 7 (cố định ở mọi lần chạy)\\
\bottomrule
\end{tabular}
\end{table}

Tổng cộng panel có $20 \times 108 = 2{.}160$ yêu cầu. Toàn bộ thực nghiệm dùng chung một hạt
giống nên hoàn toàn tất định: chạy lại cùng cấu hình cho ra cùng kết quả tới từng bit, điều này
đã được xác nhận qua các cấu hình đối chứng.

\subsection{Quy mô thực nghiệm}

Bảng \ref{bang:quymo} tổng hợp toàn bộ khối lượng thực nghiệm đã thực hiện.

\begin{table}[htbp]
\centering
\caption{Tổng hợp các đợt thực nghiệm}
\label{bang:quymo}
\small
\begin{tabular}{@{}lrrp{5.6cm}@{}}
\toprule
\textbf{Đợt} & \textbf{Lần chạy} & \textbf{Dung lượng} & \textbf{Mục đích}\\
\midrule
Panel niêm phong & 60 & 3{,}39 GB & Hai cấu hình gốc B1 và C1\\
Đường biên & 141 & 7{,}77 GB & Chín cấu hình cam kết\\
Cơ sở kết quả & 44 & 2{,}58 GB & So sánh hai cơ sở tính\\
Thí nghiệm ratchet & 80 & 4{,}30 GB & Bảo đảm một chiều, bộ chạy riêng\\
Trần cửa sổ đón & 40 & 3{,}06 GB & Nới cửa sổ 600 lên 900 giây\\
Trôi dạt (cửa sổ 600 s) & 46 & 2{,}58 GB & Bơm tắc đường đo thật, $k \in \{0,1,2\}$\\
Trôi dạt (cửa sổ 900 s) & 15 & 1{,}10 GB & Biến thể của đợt trên\\
Trôi dạt mức phân số & 12 & 1{,}22 GB & Mức trung gian $k = 0{,}5$\\
Mặt cong hai chiều & 80 & 6{,}03 GB & Thang ngân sách ở cửa sổ 900 s\\
Đặt trước ngân sách & 20 & 1{,}20 GB & Tầng cảnh báo trong hàm mục tiêu\\
Bộ sinh biết sắp xếp lại & 20 & 0{,}39 GB & Sửa chữa cùng xe, hạn mức bốn khách\\
\midrule
\textbf{Tổng} & \textbf{558} & \textbf{33{,}62 GB} & \\
\bottomrule
\end{tabular}
\end{table}

\subsection{Các cấu hình được so sánh}

Cấu hình \textbf{B1} là đường cơ sở không cam kết: hệ thống tối ưu tự do, không có sổ cái, không
có trần. Cần lưu ý B1 \emph{không phải} là một đường cơ sở cho phép tái gán tự do --- việc
chuyển một hành khách đã nhận sang xe khác bị cấm ở mọi cấu hình bởi một ràng buộc vật lý riêng.
Nghĩa là con số quy trách ở mục \ref{muc:quytrach} chỉ tính phần do \emph{đổi lịch trình}, chưa tính phần do
\emph{đổi xe} mà một hệ thống cho phép tái gán sẽ còn cộng thêm.

Cấu hình \textbf{C1} là cấu hình cam kết đầy đủ: trần 30 giây cho tổng độ dời giờ trả, khóa cứng
xe và hai điểm dừng, giờ đón cố định từ lúc xác nhận. Bảy cấu hình còn lại nằm giữa hai cực này
và tạo thành hai thang đo được mô tả ở mục \ref{muc:duongbien}.

\section{Quy trách: ai gây ra xáo trộn}
\label{muc:quytrach}

\subsection{Con số trung tâm}

Câu hỏi đầu tiên và cơ bản nhất: trong toàn bộ lượng dịch chuyển giờ hẹn mà hành khách nhìn
thấy, bao nhiêu phần là do bộ điều phối tự gây ra, bao nhiêu phần là do thế giới vật lý?

Kết quả, đo trên toàn bộ 40 ô của hai cấu hình gốc bằng cách đọc lại biên bản đầy đủ:

\begin{table}[htbp]
\centering
\caption{Phân rã lượng dịch chuyển giờ hẹn theo nguyên nhân}
\label{bang:quytrach}
\small
\begin{tabular}{@{}lrrr@{}}
\toprule
\textbf{Đại lượng} & \textbf{B1} & \textbf{C1} & \textbf{Chênh}\\
\midrule
Hành khách được phục vụ & 1.735 & 1.581 & $-154$\\
Tỉ lệ phục vụ & $80{,}32\%$ & $73{,}19\%$ & $-7{,}13$ pp\\
Số lần công bố lời hứa & 54.311 & 45.542 & $-8.769$\\
\midrule
Dịch chuyển khách nhìn thấy & 83.587 s & 9.332 s & $-74.255$ s\\
\quad --- do quyết định của nhà xe & 74.443 s & 128 s & $-74.315$ s\\
\quad --- do trôi dạt vật lý & 9.203 s & 9.204 s & $+1$ s\\
\midrule
\textbf{Tỉ lệ do nhà xe gây ra} & \textbf{89{,}06\%} & \textbf{1{,}37\%} & \\
\bottomrule
\end{tabular}
\end{table}

Con số $89{,}06\%$ là kết quả chính của mục này. Nó nói rằng trong một hệ ride-pooling tối ưu
liên tục theo lối kinh điển, gần chín phần mười sự bất định mà hành khách trải nghiệm về giờ
giấc \emph{là do hệ thống tự tạo ra}, không phải do thế giới vật lý. Và vì bộ sinh phương án
chỉ chèn thêm, toàn bộ 74.443 giây đó là dịch chuyển theo hướng xấu đi --- không có lần nào
tốt lên.

\subsection{Nhóm chứng tự nhiên}

Điều làm con số trên đáng tin là sự tồn tại của một nhóm chứng nội tại. Phần trôi dạt vật lý
giữa hai cấu hình gần như bằng nhau: 9.203 so với 9.204 giây. Thế giới vật lý cư xử như nhau ở
cả hai nhánh, còn khác biệt nằm trọn trong hành vi của bộ điều phối.

Tuy vậy, cần phát biểu điều này cho đúng. Khi mở rộng sang cấu hình thứ ba, độ lệch của thành
phần ngoại sinh lên tới $5{,}4\%$, nên việc hai con số của riêng cặp B1/C1 khớp nhau tới
$0{,}02\%$ là may mắn chứ không phải quy luật. Phát biểu chính xác là:

\begin{khungy}
Qua sáu cấu hình rất khác nhau, trôi dạt vật lý dao động trong biên độ $5{,}6\%$ và không có xu
hướng, trong khi xáo trộn do nhà điều hành gây ra thay đổi 580 lần. Hai đại lượng phản ứng hoàn
toàn khác nhau với chính sách --- đó chính là điều phép phân rã ba mốc cần chứng minh.
\end{khungy}

\subsection{Thiệt hại tập trung chứ không dàn đều}

Bảng \ref{bang:taptrung} cho thấy một đặc điểm quan trọng: thiệt hại không dàn đều mà dồn vào
một cái đuôi mỏng.

\begin{table}[htbp]
\centering
\caption{Phân bố thiệt hại trên hành khách}
\label{bang:taptrung}
\small
\begin{tabular}{@{}lrr@{}}
\toprule
& \textbf{B1} & \textbf{C1}\\
\midrule
Hành khách được phục vụ & 1.735 & 1.581\\
Khách từng bị dời do nhà xe & 403 ($23{,}2\%$) & 11 ($0{,}7\%$)\\
Tổng số lần dời & 479 & 11\\
Trung bình mỗi lần dời & 155 s & 11{,}6 s\\
Hành khách chịu thiệt nặng nhất & \textbf{840{,}9 s} & 27{,}6 s\\
Phân vị 99 & 536{,}4 s & 0{,}0 s\\
Phân vị 95 & 264{,}6 s & 0{,}0 s\\
\bottomrule
\end{tabular}
\end{table}

Ba điều đọc được từ bảng này. Thứ nhất, ba phần tư hành khách không hề bị ảnh hưởng --- vấn đề
nằm ở một cái đuôi mỏng. Thứ hai, mỗi lần bộ điều phối đổi ý, nó dời giờ đã báo trung bình 155
giây; đây không phải chuyện nhích vài giây mà mỗi lần là hơn hai phút rưỡi. Thứ ba, hành khách
chịu thiệt nặng nhất bị dời tổng cộng gần mười bốn phút mà hệ thống vẫn tính là phục vụ thành
công, vì mọi mốc vẫn nằm trong cửa sổ $\Omega$ và $\Delta$ hợp lệ.

Điểm thứ ba là minh họa cụ thể nhất cho hạn chế thứ nhất đã nêu ở mục 1.2, và nó cũng cho thấy
vì sao một số hạng cộng gộp trong hàm mục tiêu là không đủ: khi cộng gộp, mười bốn phút của một
người bị pha loãng trong 1.735 người.

\subsection{Độ khít của phép phân rã}

Vì cả ba đại lượng đều là tổng trị tuyệt đối, bất đẳng thức tam giác chỉ cho $B \le A + d$.
Bảng \ref{bang:cong-khit} kiểm tra độ hụt thực tế.

\begin{table}[htbp]
\centering
\caption{Độ khít của phân rã $B \le A + d$ trên sáu cấu hình (đơn vị: giây)}
\label{bang:cong-khit}
\small
\begin{tabular}{@{}lrrrrrr@{}}
\toprule
\textbf{Cấu hình} & $A$ & $d$ & $A+d$ & $B$ & \textbf{hụt} & \textbf{hụt \%}\\
\midrule
B1 & 74.443 & 9.203 & 83.646 & 83.587 & 58 & $0{,}07\%$\\
nobudget & 22.485 & 9.703 & 32.188 & 32.176 & 12 & $0{,}04\%$\\
budget300 & 15.809 & 9.641 & 25.450 & 25.441 & 9 & $0{,}03\%$\\
budget120 & 5.063 & 9.260 & 14.324 & 14.318 & 5 & $0{,}04\%$\\
budget60 & 1.109 & 9.184 & 10.293 & 10.291 & 1 & $0{,}01\%$\\
C1 & 128 & 9.204 & 9.332 & 9.332 & 0 & $\mathbf{0{,}00\%}$\\
\bottomrule
\end{tabular}
\end{table}

Độ hụt tối đa là $0{,}07\%$ và bằng đúng không ở C1 sau khi làm tròn.

Con số gộp này dễ bị đọc quá mạnh, nên mục \ref{muc:phanra} đã kiểm tra lại ở mức \emph{từng lần
công bố} thay vì trên tổng (bảng \ref{bang:khit-tung-lan}). Kết quả mịn hơn và khác về chất: dấu
bằng đúng ở $99{,}58\%$ số lần công bố, nhưng phần lớn con số đó đến từ việc \emph{một trong hai}
thành phần bằng không --- điều xảy ra ở $99{,}12\%$ số lần, khiến phép kiểm trở nên tầm thường.
Giới hạn ở 476 lần công bố mà cả $A$ và $d$ đều khác không, dấu bằng chỉ đúng $52{,}1\%$; trong
228 lần còn lại, trôi dạt \emph{thuận chiều} đã hấp thụ một phần cú đẩy của bộ điều phối, cho
$B = A - d$ thay vì $B = A + d$. Độ hụt lớn nhất trong một lần công bố là $1{,}38$ giây.

Tổng phần hụt vẫn chỉ là $0{,}0748\%$ trên nhánh B1, nhỏ hơn sàn nhiễu quỹ đạo hai bậc độ lớn,
nên không kết luận định lượng nào trong chương này thay đổi. Nhưng cách phát biểu thì phải đổi:
đây là một \emph{bất đẳng thức có độ khít đo được}, không phải một đẳng thức.

Một lưu ý về cách phát biểu: vì phép phân rã không phải một phân hoạch chính xác, không nên nói
``$89{,}06\%$ lượng xáo trộn \emph{được gây ra bởi} bộ điều phối'' mà nên nói ``phần do bộ điều
phối gây ra \emph{bằng} $89{,}06\%$ lượng khách nhìn thấy'' --- một tỉ số, đúng theo nghĩa đen.
Sai khác chỉ $0{,}07\%$ nên kết luận không đổi, nhưng cách nói cần chính xác.

\section{Đường biên: giá của cam kết}
\label{muc:duongbien}

\subsection{Thiết kế hai thang}

Để đo giá của cam kết một cách có kiểm soát, chín cấu hình được xếp thành hai thang cắt nhau
tại C1. Thang thứ nhất giữ nguyên mọi khóa và chỉ thay đổi trần ngân sách giờ trả. Thang thứ
hai giữ trần ở 30 giây và thay đổi cách bảo vệ giờ đón.

\subsection{Thang ngân sách}

\begin{table}[htbp]
\centering
\caption{Thang ngân sách giờ trả --- đường biên giữa phục vụ và ổn định}
\label{bang:thang1}
\small
\begin{tabular}{@{}llrrrrr@{}}
\toprule
\textbf{Cấu hình} & \textbf{Lời hứa} & \textbf{Phục vụ} & \textbf{vs B1} & \textbf{Xáo trộn} & \textbf{Tệ nhất} & \textbf{Bị dời}\\
\midrule
B1 & không hứa gì & 1.735 & 0 & 74.443 s & 841 s & 403\\
nobudget & khóa xe + điểm & 1.655 & $-3{,}70$ & 22.485 s & 595 s & 142\\
budget300 & $+$ trần 300 s & 1.641 & $-4{,}35$ & 15.809 s & 291 s & 128\\
budget120 & $+$ trần 120 s & 1.608 & $-5{,}88$ & 5.063 s & 119 s & 76\\
budget60 & $+$ trần 60 s & 1.579 & $-7{,}22$ & 1.109 s & 59 s & 34\\
C1 & $+$ trần 30 s & 1.581 & $-7{,}13$ & 128 s & 28 s & 11\\
\bottomrule
\end{tabular}
\end{table}

Đây là bảng trung tâm của khóa luận. Đọc từ trên xuống, nó cho thấy chính xác cái giá của từng
mức chặt chẽ. Bốn nhận xét đáng chú ý.

\textbf{Trần luôn được tôn trọng.} Ở mỗi cấu hình, hành khách chịu thiệt nặng nhất đều nằm dưới
trần: 291 dưới 300, 119 dưới 120, 59 dưới 60, và 27{,}6 dưới 30. Đây là bằng chứng rằng lời hứa
cưỡng chế được chứ không chỉ là chỉ tiêu mong muốn.

\textbf{Phần đắt nhất là quyết định cam kết chút nào đó.} Bước từ B1 sang nobudget --- chỉ khóa
danh tính xe và hai điểm dừng, chưa hứa một chữ nào về thời gian --- đã ngốn $3{,}70$ điểm phần
trăm, tức hơn nửa tổng chi phí $7{,}13$ pp.

\textbf{Điểm gãy nằm ở 300 giây.} Với $4{,}35$ pp, hệ thống mua được lời hứa ``giờ trả không xê
dịch quá năm phút'' --- chỉ thêm $0{,}65$ pp so với việc chỉ khóa danh tính, mà cắt đôi trường
hợp tệ nhất từ 595 xuống 291 giây.

\textbf{Hiệu suất mua sự ổn định giảm dần.} Số giây ổn định mua được trên mỗi chuyến phải hy
sinh lần lượt là 649, 624, 546, 470, 483 --- giảm đều cho tới khi chạm sàn nhiễu ở hai điểm cuối.

\subsection{Cảnh báo phương pháp luận: độ nhạy quỹ đạo}

Bảng \ref{bang:thang1} chứa một điều bất thường cần giải thích: cấu hình budget60 phục vụ 1.579
hành khách còn C1 với trần chặt hơn lại phục vụ 1.581. Trần chặt hơn cho kết quả tốt hơn.

Kiểm tra từng ô cho thấy 11 trên 20 ô khác nhau giữa hai cấu hình, và khác theo \emph{cả hai
chiều}; 8 trên 20 ô vi phạm thứ tự đơn điệu. Đây không phải nhiễu thống kê --- mọi lần chạy dùng
cùng hạt giống và hoàn toàn tất định. Đây là \emph{độ nhạy quỹ đạo}: ràng buộc chặt hơn làm thay
đổi phương án nào được chọn ở những bước sớm, và khác biệt đó lan truyền qua chuỗi quyết định
tiếp theo. Tất định nhưng thực tế không dự đoán được, theo cùng cơ chế khiến một thuật toán tham
lam không cho nghiệm tối ưu toàn cục.

Hệ quả bắt buộc là phải công bố một sàn nhiễu. Độ tán ở mức ô là $\pm 1$ đến $\pm 4$ hành khách
trên khoảng 80, nên trên bảng 20 ô, \emph{chênh lệch dưới khoảng $\pm 7$ hành khách, tức
$\pm 0{,}3$ điểm phần trăm, không diễn giải được}.

\begin{table}[htbp]
\centering
\caption{Áp dụng sàn nhiễu cho các bước trên thang ngân sách}
\label{bang:sannhieu}
\small
\begin{tabular}{@{}lrl@{}}
\toprule
\textbf{Bước} & \textbf{Chênh lệch} & \textbf{Kết luận}\\
\midrule
$\infty \to 300$ s & $0{,}65$ pp (14 khách) & vừa đủ trên sàn\\
$300 \to 120$ s & $1{,}53$ pp (33 khách) & kết quả thật\\
$120 \to 60$ s & $1{,}34$ pp (29 khách) & kết quả thật\\
$60 \to 30$ s & $0{,}09$ pp (2 khách) & \textbf{thuần nhiễu quỹ đạo}\\
\bottomrule
\end{tabular}
\end{table}

Không có bước kiểm tra từng ô này, báo cáo đã có thể kết luận rằng ``siết từ 60 xuống 30 giây
tiết kiệm được hai chuyến'' --- một kết luận vô nghĩa. Đây là một trong những bài học phương
pháp quan trọng nhất rút ra được trong quá trình thực hiện.

\subsection{Thang bảo vệ giờ đón}

\begin{table}[htbp]
\centering
\caption{Thang bảo vệ giờ đón --- sáu cơ chế rơi vào hai hành vi}
\label{bang:thang2}
\small
\begin{tabular}{@{}llrrrr@{}}
\toprule
\textbf{Cấu hình} & \textbf{Cách bảo vệ giờ đón} & \textbf{Phục vụ} & \textbf{vs B1} & \textbf{Xáo trộn} & \textbf{Tệ nhất}\\
\midrule
nopickuplock & không bảo vệ & 1.583 & $-7{,}04$ & 187 s & 34 s\\
freeze300 & hàng rào 300 s & 1.583 & $-7{,}04$ & 187 s & 34 s\\
freeze600 & hàng rào 600 s & 1.581 & $-7{,}13$ & 128 s & 28 s\\
C1 & khóa khi xác nhận & 1.581 & $-7{,}13$ & 128 s & 28 s\\
\bottomrule
\end{tabular}
\end{table}

Thang này cho một kết quả bất ngờ: bốn cơ chế khác nhau chỉ sinh ra \emph{hai} bộ số. Lý do được
làm rõ bằng một phép đo bổ sung. Quét 19.539 lần công bố lời hứa và tính khoảng cách từ lúc công
bố tới lúc đón:

\begin{table}[htbp]
\centering
\caption{Thời điểm các lần công bố lời hứa so với giờ đón}
\label{bang:thoidiem}
\small
\begin{tabular}{@{}lrr@{}}
\toprule
\textbf{Thời điểm công bố} & \textbf{Số lần} & \textbf{Tỉ lệ}\\
\midrule
Khách đã lên xe & 9.057 & $46{,}4\%$\\
Còn 0--300 giây & 7.199 & $36{,}8\%$\\
Còn 300--600 giây & 3.283 & $16{,}8\%$\\
Còn hơn 600 giây & \textbf{0} & $\mathbf{0{,}0\%}$\\
\bottomrule
\end{tabular}
\end{table}

Ba hệ quả. Thứ nhất, \emph{hàng rào 600 giây không bao giờ mở}, vì tầm đặt chỗ của kịch bản đúng
bằng 600 giây; mọi hàng rào từ 600 giây trở lên đều tương đương khóa vĩnh viễn. Đây là một giới
hạn phạm vi phải ghi rõ: khóa luận không kết luận được gì về hàng rào dài hơn tầm đặt chỗ. Thứ
hai, \emph{hàng rào 300 giây tương đương không khóa}, vì trong 300 giây cuối xe đã cam kết và
đang trên đường, giờ đón vốn không còn chỗ xê dịch. Thứ ba, $46{,}4\%$ số lần công bố xảy ra sau
khi khách đã lên xe nên chúng chỉ nói về giờ trả --- đây là lời giải thích cơ học cho sự bất đối
xứng giữa hai thang: ngân sách giờ trả chi phối gần một nửa khối lượng công bố nên đắt, còn khóa
giờ đón chỉ chạm vào phần đầu vốn đã chật nên rẻ.

Ba kết quả độc lập --- chi phí thang một, chi phí thang hai, và phân bố thời điểm công bố ---
khớp vào cùng một lời giải thích.

\section{So sánh hai cơ sở tính}
\label{muc:haicoso}

Mục này kiểm chứng thực nghiệm ba định lý ở mục 3.5.

\subsection{Kết quả ở cùng mức ngân sách}

\begin{table}[htbp]
\centering
\caption{Cơ sở hành vi so với cơ sở kết quả ở cùng trần}
\label{bang:coso}
\small
\begin{tabular}{@{}lrrl@{}}
\toprule
\textbf{Trần} & \textbf{Hành vi khả thi} & \textbf{Kết quả khả thi} & \textbf{Phục vụ trên các ô chung}\\
\midrule
30 s & \textbf{20/20} & 16/20 & hành vi 1.294 --- kết quả 1.294 --- chênh \textbf{0}\\
120 s & \textbf{20/20} & 16/20 & hành vi 1.267 --- kết quả 1.266 --- chênh $-1$\\
\bottomrule
\end{tabular}
\end{table}

Kết quả khớp chính xác với lý thuyết. Định lý 3.1 dự đoán cơ sở hành vi không bao giờ bị làm
rỗng --- quan sát được 20/20 ô ở cả hai trần. Định lý 3.2 dự đoán cơ sở kết quả có thể vô nghiệm
--- quan sát được 4/20 ô hỏng, với đúng thông báo lỗi ``đã loại mọi phương án sinh ra cho xe
này''. Định lý 3.3 dự đoán cơ sở kết quả không phục vụ nhiều hơn --- quan sát được chênh lệch 0
và $-1$ hành khách.

\subsection{Nới trần không chữa được, chỉ dời chỗ đau}

Một câu hỏi tự nhiên là liệu nới trần có làm cơ sở kết quả hết vô nghiệm không. Bảng
\ref{bang:doicho} cho câu trả lời.

\begin{table}[htbp]
\centering
\caption{Tập ô vô nghiệm dưới cơ sở kết quả, ở hai mức trần}
\label{bang:doicho}
\small
\begin{tabular}{@{}ll@{}}
\toprule
\textbf{Trần} & \textbf{Các ô vô nghiệm}\\
\midrule
30 s & 16-r2, \textbf{16-r4}, 17-r1, 17-r3\\
120 s & 14-r4, 15-r2, \textbf{16-r4}, 17-r4\\
\bottomrule
\end{tabular}
\end{table}

Nới trần gấp bốn lần vẫn hỏng đúng bốn ô, nhưng là một \emph{tập ô khác}, giao nhau ở đúng một
ô. Đó là sự dời chỗ chứ không phải sự chữa lành.

\subsection{Diễn giải}

Kết hợp lý thuyết và thực nghiệm cho một kết luận thiết kế dứt khoát. Lý thuyết nói cơ sở hành
vi luôn khả thi, cơ sở kết quả có thể vô nghiệm, và cơ sở kết quả không bao giờ hơn. Thực nghiệm
nói thêm một điều mà lý thuyết không nói: \emph{khi cơ sở kết quả khả thi thì nó cũng chẳng cho
thêm gì} --- chênh 0 và $-1$ hành khách.

\begin{khungy}
Nếu muốn hứa với hành khách một cách cưỡng chế được, phải hứa về \emph{hành vi của mình} ---
``tôi sẽ không tự dời giờ của bạn quá $\beta$'' --- chứ không thể hứa về \emph{kết quả} ---
``giờ của bạn sẽ không đổi quá $\beta$''. Lời hứa thứ hai nghe hay hơn nhưng không giữ được, vì
một phần của nó nằm ngoài tầm kiểm soát của người hứa. Và nó không mua thêm được gì cả.
\end{khungy}

\section{Độ vững với cách đo}
\label{muc:dovung}

\subsection{Vì sao phải kiểm tra}

RideBound mặc định dùng một thành viên của họ thước đo bất ổn: tổng tuyến tính không trọng số.
Như mục 2.3 đã nêu, văn liệu hoạch định sản xuất cung cấp cả một họ thước đo cho phép đặt trọng
số khác nhau theo vị trí trên chân trời kế hoạch và dùng cả biểu thức bậc hai. Nếu thứ hạng các
cấu hình đổi khi đổi thước đo thì chính điều đó là một phát hiện.

\subsection{Kết quả}

\begin{table}[htbp]
\centering
\caption{Chín thước đo áp dụng song song trên sáu cấu hình}
\label{bang:chinthuoc}
\footnotesize
\begin{tabular}{@{}lrrrrrr@{}}
\toprule
\textbf{Thước đo} & \textbf{B1} & \textbf{nobudget} & \textbf{300 s} & \textbf{120 s} & \textbf{60 s} & \textbf{30 s}\\
\midrule
Tuyến tính $\sum|\delta|$ & 74.443 & 22.485 & 15.809 & 5.063 & 1.109 & 128\\
Bậc hai $\sqrt{\sum \delta^2}$ & 4.502 & 2.288 & 1.590 & 649 & 227 & 52\\
Trọng số bậc thang & 180.824 & 58.841 & 40.943 & 13.129 & 2.975 & 300\\
Trọng số cận kề & 41.571 & 13.618 & 9.541 & 3.214 & 728 & 77\\
Đếm số lần dời & 479 & 155 & 135 & 79 & 34 & 11\\
Khách tệ nhất & 841 & 595 & 291 & 119 & 59 & 28\\
Phân vị 99 & 536 & 331 & 218 & 101 & 39 & 0\\
Phân vị 95 & 265 & 115 & 83 & 0 & 0 & 0\\
Số khách bị dời & 403 & 142 & 128 & 76 & 34 & 11\\
\bottomrule
\end{tabular}
\end{table}

Cả chín thước đo đơn điệu trên cả sáu cấu hình, không một ngoại lệ. Kết luận chính của khóa luận
không phụ thuộc vào lựa chọn cách đo.

Tuy nhiên, \emph{độ lớn} thì lệch tới khoảng 300 lần. Thước tuyến tính nói C1 tốt hơn B1 khoảng
585 lần; thước đếm số lần nói 43 lần; thước khách-tệ-nhất nói 30 lần. Cùng dữ liệu, cùng hướng,
hai bậc độ lớn khác nhau. Hệ quả thực hành là không được quảng cáo con số 585 lần; phải báo cáo
cả họ thước đo và nói rõ con số nào ứng với câu hỏi nào.

\subsection{Bất đối xứng giữa ba đại lượng}

Phân tích trên dẫn tới một quan sát mà chúng tôi coi là kết quả phương pháp trung tâm.

\begin{table}[htbp]
\centering
\caption{Ba đại lượng phản ứng khác hẳn nhau với chính sách}
\label{bang:batdoixung}
\small
\begin{tabular}{@{}lp{9.4cm}@{}}
\toprule
\textbf{Đại lượng} & \textbf{Phản ứng với chính sách}\\
\midrule
Độ ổn định lời hứa & Đơn điệu tuyệt đối, 9 thước đo, 0 ngoại lệ, trải 175 lần.\\
Trôi dạt vật lý & Gần như bất biến, biên độ $\pm 5{,}6\%$, không có xu hướng.\\
Tỉ lệ phục vụ & Nhạy quỹ đạo, vi phạm đơn điệu ở 8/20 ô, sàn nhiễu $\pm 0{,}3$ pp.\\
\bottomrule
\end{tabular}
\end{table}

Phép phân rã ba mốc tách được một đại lượng \emph{phản ứng sạch với chính sách} khỏi một đại
lượng \emph{không phản ứng}, trong khi thước đo truyền thống --- tỉ lệ phục vụ --- lại là thứ
nhiễu nhất trong ba. Một báo cáo chỉ dựa vào tỉ lệ phục vụ rất dễ đọc nhầm nhiễu thành kết quả.

\section{Thí nghiệm phụ: bảo đảm một chiều}
\label{muc:ratchet}

\subsection{Vì sao phải tách riêng}

Hai cấu hình dùng cơ chế bảo đảm một chiều không chạy được trên bộ chạy đã niêm phong. Nguyên
nhân xác định được: bộ chạy niêm phong được dựng ngày 22 tháng 8 năm 2026, còn tính năng ratchet
vào mã nguồn ngày 4 tháng 9 --- sau mười ba ngày. Bộ chạy không biết trường cấu hình đó nên
thoát ngay ở bước bắt tay.

Vì kết quả được ghim theo băm nhị phân bộ chạy, không thể chỉ đơn giản dựng bộ chạy mới rồi trộn
số liệu vào bảng chính. Giải pháp là một thí nghiệm phụ tách bạch: dựng bộ chạy mới trong thư
mục riêng, chạy cả hai cấu hình điều trị \emph{và} hai cấu hình đối chứng đã biết kết quả trên
bộ chạy cũ. Nếu hai đối chứng tái lập đúng, số liệu điều trị đọc được trên cùng thang.

\subsection{Kiểm chứng bộ chạy mới}

\begin{table}[htbp]
\centering
\caption{Hai cấu hình đối chứng tái lập kết quả của bộ chạy niêm phong}
\label{bang:doichung}
\small
\begin{tabular}{@{}lrrr@{}}
\toprule
\textbf{Đối chứng} & \textbf{Bộ chạy mới} & \textbf{Bộ chạy niêm phong} & \textbf{Ô sai lệch}\\
\midrule
reference (khóa cứng) & 1.581 & 1.581 & \textbf{0/20}\\
nopickuplock (không khóa) & 1.583 & 1.583 & \textbf{0/20}\\
\bottomrule
\end{tabular}
\end{table}

Hai cấu hình ở hai cực của thang, tái lập từng ô một. Bộ chạy mới tương đương về hành vi.

\subsection{Kết quả và lời giải thích}

\begin{table}[htbp]
\centering
\caption{Bảo đảm một chiều so với khóa cứng hai chiều}
\label{bang:ratchet}
\small
\begin{tabular}{@{}lrrr@{}}
\toprule
& \textbf{nopickuplock} & \textbf{ratchet} & \textbf{khóa cứng}\\
\midrule
Phục vụ & 1.583 & \textbf{1.581} & 1.581\\
Xáo trộn do quyết định & 187 s & \textbf{128 s} & 128 s\\
Số lần dời & 13 & \textbf{11} & 11\\
Khách tệ nhất & 34{,}4 s & \textbf{27{,}6 s} & 27{,}6 s\\
\bottomrule
\end{tabular}
\end{table}

Bảo đảm một chiều cho kết quả \emph{giống hệt} khóa cứng hai chiều, không lệch một chữ số. Lời
giải thích có thể rút ra chỉ từ ba con số trong bảng: cấu hình thả tự do cho 59 giây xê dịch giờ
đón do quyết định ($187 - 128$), cấu hình ratchet cho 0 giây, vậy \emph{toàn bộ 59 giây đó đều
là xấu đi}. Nếu có phần nào là cải thiện, ratchet đã cho qua.

Nguyên nhân mang tính cơ học: bộ sinh phương án chỉ chèn thêm, không bỏ bớt, nên chỉ có thể đẩy
lùi giờ đón của khách cũ chứ không bao giờ kéo sớm lên. Quyền ``được phép cải thiện'' mà ratchet
trao là một quyền không bao giờ dùng được. Điều này khớp với kết quả đã biết
$\mathtt{decisionMovedEarlier} = 0$ trên 421.752 lần sửa.

Kết quả này có ý nghĩa đối chiếu văn liệu trực tiếp, vì ratchet chính là cơ chế Alonso-Mora và
cộng sự sử dụng:

\begin{khungy}
Trong một bộ điều phối dựa trên chèn, bảo đảm một chiều nghiêm ngặt \emph{đúng bằng} khóa cứng.
Nó trông mềm hơn nhưng không mềm hơn chút nào. Hình thức của một lời bảo đảm chỉ có ý nghĩa
trong chừng mực bộ sinh phương án có thể thực thi nó; một bảo đảm nới lỏng theo hướng mà thuật
toán không bao giờ đi được thì không nới lỏng gì cả.
\end{khungy}

Cuối cùng, cần ghi nhận vai trò của cơ chế chứng cứ. Nếu chỉ nhìn bảng tổng hợp, việc sáu cấu
hình rơi vào hai bộ số trông hệt như một lỗi sao chép cấu hình. Chỉ nhờ băm ngữ nghĩa và tóm
lược công bố khác nhau ở toàn bộ 20/20 ô mới biết đó là những lần chạy thật, đi qua những chuỗi
lời hứa khác nhau, và về đích cùng một chỗ.

\section{Thí nghiệm trôi dạt thật}
\label{muc:troidat}

\subsection{Động cơ}

Ba định lý ở mục 3.5 đúng nhưng gần như rỗng trên panel chính, vì một lý do được phát hiện muộn
trong quá trình thực hiện: bộ dữ liệu thử có \emph{đúng một ảnh chụp thời gian di chuyển} cho cả
hai giờ mô phỏng. Kiểm tra trực tiếp xác nhận điều này ở ba tầng --- trường
\texttt{travelSnapshots} có độ dài 1, toàn bộ 108 sự kiện đều thuộc loại ``yêu cầu xuất hiện'',
và quét toàn bộ biên bản một lần chạy chỉ thấy một định danh ảnh chụp duy nhất.

Nghĩa là trôi dạt $d$ đo được trên panel chính \emph{không đến từ tắc đường} mà đến từ độ trễ
thực thi: xe di chuyển theo bước rời rạc và có thời gian dừng đón trả, nên khi chiếu lại cùng
một lộ trình từ vị trí thật hiện tại thì giờ đến lệch đi một chút. Đó là lý do $d$ nhỏ --- trung
vị $4{,}9$ giây, tối đa $39{,}9$ giây. Đây là \emph{cận dưới} của trôi dạt, không phải giá trị
điển hình của một hệ thật.

Thí nghiệm này bơm trôi dạt thật vào để kiểm chứng ba định lý trong chế độ mà chúng có ý nghĩa.

\subsection{Cách bơm trôi dạt}

Trôi dạt không phải nhiễu nhân tạo. Bộ dữ liệu FleetPy Manhattan có sẵn hệ số thời gian di
chuyển \emph{đo thật}, mười lăm phút một mốc, riêng cho từng ngày. Trong khung 08:00--10:00 của
panel, cao điểm sáng làm thời gian di chuyển tăng từ $17\%$ đến $26\%$ tùy ngày, và tăng đơn
điệu. Hệ số được áp dụng theo công thức

\begin{equation}
f_k(t) = \left( \frac{f_{\text{thật}}(t_0 + t)}{f_{\text{thật}}(t_0)} \right)^{k}
\end{equation}

trong đó $k$ là mức khuếch đại: $k=0$ tái lập chính xác đường cơ sở phẳng, $k=1$ là ngày thật,
$k=2$ nhân đôi độ lệch. Việc chuẩn hóa theo mốc đầu khung bảo đảm $f_k(0) = 1$ tuyệt đối, điều
này cần thiết vì hệ thống kiểm tra ở bước khởi tạo rằng bộ định tuyến tái tạo đúng thời gian di
chuyển của bộ dữ liệu thử; trôi dạt chỉ tích lũy \emph{sau} phép kiểm tra đó.

\subsection{Trôi dạt đo được}

\begin{table}[htbp]
\centering
\caption{Trôi dạt tích lũy trên mỗi hành khách ở ba mức}
\label{bang:troidat}
\small
\begin{tabular}{@{}lrrrr@{}}
\toprule
\textbf{Mức} & $d$ \textbf{p50} & $d$ \textbf{p90} & $d$ \textbf{max} & \textbf{\% khách vượt trần 30 s vì riêng $d$}\\
\midrule
$k=0$ phẳng & $4{,}9$ s & $10{,}8$ s & $39{,}9$ s & $0{,}2\%$\\
$k=1$ Manhattan thật & $14{,}1$ s & $47{,}4$ s & $150{,}7$ s & $\mathbf{22{,}8\%}$\\
$k=2$ nhân đôi & $10{,}9$ s & $69{,}9$ s & $157{,}0$ s & $20{,}3\%$\\
\bottomrule
\end{tabular}
\end{table}

Con số cột cuối là mấu chốt: dưới tắc đường Manhattan thật, gần một phần tư số hành khách tích
lũy hơn 30 giây trôi dạt \emph{thuần túy do giao thông}. Dưới cơ sở kết quả, ngân sách của những
người đó bị thủng vì thời tiết --- nhà điều hành bị tính phí cho việc không làm gì.

Cần ghi kèm một lưu ý: số liệu $k=1$ và $k=2$ chỉ tính trên những ô sống sót qua bức tường mô tả
ở mục \ref{muc:batbien}, nên có thiên lệch sống sót. Con số $k=0$ là trên đủ 20 ô.

\subsection{Kết quả chính}

Điều quan trọng nhất trong thí nghiệm này là việc phân loại nguyên nhân hỏng. Đếm số lần hỏng
thôi là không đủ --- phải mở biên bản lỗi của từng lần chạy.

\begin{table}[htbp]
\centering
\caption{Phân loại kết quả theo nguyên nhân thật}
\label{bang:phanloai}
\small
\begin{tabular}{@{}lrrr@{}}
\toprule
\textbf{Cấu hình} & \textbf{Chạy được} & \textbf{Vô nghiệm do cam kết} & \textbf{Hỏng do ràng buộc vật lý}\\
\midrule
hành vi, $k=0$ & 20 & \textbf{0} & 0\\
kết quả, $k=0$ & 16 & \textbf{4} & 0\\
hành vi, $k=1$ & 9 & \textbf{0} & 11\\
kết quả, $k=1$ & 0 & \textbf{19} & 1\\
\bottomrule
\end{tabular}
\end{table}

Kết quả xác nhận cả ba định lý trong chế độ có trôi dạt thật.

\textbf{Định lý 3.1 đứng vững.} Cơ sở hành vi có đúng \emph{không} lần vô nghiệm do cam kết ở
mọi mức trôi dạt. Mười một lần hỏng ở $k=1$ thuộc một loại hoàn toàn khác --- một bất biến miền
của mô hình, mô tả ở mục \ref{muc:batbien} --- và loại lỗi ấy đập vào cả hai cơ sở như nhau. Nếu chỉ đếm số
lần hỏng mà không mở biên bản, kết luận rút ra sẽ là ``định lý sai'', ngược hoàn toàn.

\textbf{Định lý 3.2 được xác nhận và mạnh hơn dự kiến.} Cơ sở kết quả vô nghiệm 4/20 ô khi trôi
dạt ở mức cận dưới, và \emph{19/20 ô} dưới tắc đường Manhattan thật. Ở mức đó, điều phối viên
không còn \emph{bất kỳ} quyết định hợp lệ nào, kể cả không làm gì.

\subsection{Quy trách dưới trôi dạt thật}

Một phép đo bổ sung, thực hiện trên cùng chín ô mà cả hai mức đều chạy được, cho thấy phần đóng
góp của nhà điều hành teo lại thế nào khi thế giới ồn hơn.

\begin{table}[htbp]
\centering
\caption{Quy trách dưới chế độ cam kết, so hai mức trôi dạt (9 ô chung)}
\label{bang:quytrach-troidat}
\small
\begin{tabular}{@{}lrrrr@{}}
\toprule
\textbf{Mức} & \textbf{Do nhà xe} & \textbf{Do đường} & \textbf{Khách thấy} & \textbf{\% do nhà xe}\\
\midrule
$k=0$ phẳng & 71 s & 3.984 s & 4.056 s & $1{,}76\%$\\
$k=1$ Manhattan thật & 27 s & 12.155 s & 12.182 s & $\mathbf{0{,}22\%}$\\
\bottomrule
\end{tabular}
\end{table}

Tỉ lệ 450 trên 1. Tính lời hứa theo kết quả nghĩa là bắt nhà điều hành chịu trách nhiệm cho
12.155 giây tắc đường để kiểm soát 27 giây hành vi của chính mình. Đây là lập luận nghĩa vụ
phương tiện ở mục 2.4, giờ đã có số: cơ sở kết quả không chỉ \emph{bất khả thi} mà còn
\emph{bất công}.

\subsection{Thang trôi dạt ba điểm}

Để kiểm tra xem kết quả trên có phải hiện tượng ngưỡng hay là một xu hướng liên tục, một mức
trung gian $k = 0{,}5$ được bổ sung. Mức này làm thời gian di chuyển tăng $8{,}3\%$ đến
$12{,}2\%$ tùy ngày, tức đúng một nửa độ lệch lô-ga-rít của ngày thật.

\begin{table}[htbp]
\centering
\caption{Thang trôi dạt ba điểm --- tính khả thi của hai cơ sở}
\label{bang:thang-troidat}
\small
\begin{tabular}{@{}lrrrrr@{}}
\toprule
\textbf{Mức} & $d$ \textbf{p50} & $d$ \textbf{p90} & \textbf{\% khách vượt 30 s} & \textbf{Hành vi} & \textbf{Kết quả}\\
\midrule
$k=0$ phẳng & $4{,}9$ s & $10{,}8$ s & $0{,}2\%$ & \textbf{0/20} & 4/20\\
$k=0{,}5$ & $9{,}8$ s & $27{,}6$ s & $7{,}9\%$ & \textbf{0/20} & 17/20\\
$k=1$ thật & $14{,}1$ s & $47{,}4$ s & $22{,}8\%$ & \textbf{0/20} & 19/20\\
\bottomrule
\end{tabular}
\end{table}

Hai cột cuối là số ô \emph{vô nghiệm do ràng buộc cam kết}. Cơ sở hành vi giữ nguyên số không ở
cả ba mức --- Định lý 3.1 đúng ở mọi mức trôi dạt. Cơ sở kết quả tăng đơn điệu từ 4 lên 17 rồi
19 ô. Đây là một thang ba điểm chứ không phải một quan sát ngưỡng.

\subsection{Vì sao không thể so tỉ lệ phục vụ --- một lập luận số học}

Thí nghiệm này cho kết quả về \emph{tính khả thi} nhưng không cho kết quả về \emph{tỉ lệ phục
vụ}, và lý do sâu hơn một trở ngại kỹ thuật. Ban đầu chúng tôi cho rằng chỉ cần tìm một mức
trôi dạt trung gian là sẽ có đủ ô mà cả hai cơ sở cùng chạy được. Dữ liệu bác bỏ điều đó.

Một ô là khả thi dưới cơ sở kết quả khi \emph{không một hành khách nào} trong ô có trôi dạt tích
lũy vượt trần. Nếu tỉ lệ hành khách vượt trần là $p$ và mỗi ô có $n$ hành khách thì, với giả
thiết độc lập xấp xỉ,

\begin{equation}
P(\text{ô khả thi}) \approx (1-p)^{\,n}, \qquad n \approx 80
\end{equation}

\begin{table}[htbp]
\centering
\caption{Mô hình khả thi cấp ô đối chiếu với quan sát}
\label{bang:mohinh-kha-thi}
\small
\begin{tabular}{@{}lrrrr@{}}
\toprule
\textbf{Mức} & $p$ & $(1-p)^{80}$ & \textbf{Dự đoán ô khả thi} & \textbf{Quan sát}\\
\midrule
$k=0$ & $0{,}2\%$ & $0{,}8520$ & $17{,}0$ & 16/20\\
$k=0{,}5$ & $7{,}9\%$ & $0{,}0014$ & $0{,}0$ & 0/17\\
$k=1$ & $22{,}8\%$ & $0{,}0000$ & $0{,}0$ & 0/19\\
\bottomrule
\end{tabular}
\end{table}

Mô hình khớp với quan sát ở cả ba mức. Đọc ngược lại, để một nửa số ô còn khả thi thì cần
$p \le 0{,}86\%$, và để $80\%$ số ô còn khả thi thì cần $p \le 0{,}28\%$ --- tức gần như đúng
mức trôi dạt của panel phẳng, nơi $p = 0{,}2\%$. Không tồn tại mức trôi dạt nào vừa đủ lớn để
hai cơ sở tách nhau, vừa đủ nhỏ để còn nhiều ô chung.

Kết luận là phép so sánh tỉ lệ phục vụ trên tập ô chung \emph{không thể lấy được về mặt cấu
trúc}, và điều đó tự nó là một phát hiện:

\begin{khungy}
Hệ quả của Định lý 3.3 biểu hiện theo kiểu \emph{phân loại} chứ không theo kiểu \emph{biên}.
Cơ sở kết quả không phục vụ ít hơn vài hành khách trên các ô chung --- nó mất nguyên cả ô. Với
khoảng 80 hành khách mỗi ô, khả thi cấp ô sụp theo hàm mũ $(1-p)^{n}$ ngay khi $p$ rời khỏi
lân cận của không.
\end{khungy}

Nói cách khác, ở mức trôi dạt nhỏ hai cơ sở gần như trùng nhau vì dải khác biệt rộng đúng $d$
mà $d$ quá nhỏ so với cú nhảy điển hình 130 giây (Hệ quả 3.4); ở mức trôi dạt thật thì cơ sở
kết quả mất trắng. \emph{Không có khoảng giữa.} Đây là câu trả lời đầy đủ cho câu hỏi ở mục
3.5.4, chỉ có điều nó không phải câu trả lời mà chúng tôi dự kiến khi thiết kế thí nghiệm.

\section{Phép thử trần cửa sổ đón và một bất biến miền}
\label{muc:cuaso}

\subsection{Động cơ}

Phân tích nguyên nhân loại bỏ phương án cho thấy cửa sổ đón chiếm ưu thế áp đảo.

\begin{table}[htbp]
\centering
\caption{Nguyên nhân loại bỏ phương án ứng viên (cấu hình B1, 20 ô)}
\label{bang:nguyennhan}
\small
\begin{tabular}{@{}lrr@{}}
\toprule
\textbf{Nguyên nhân} & \textbf{Số lần} & \textbf{Tỉ lệ}\\
\midrule
Quá cửa sổ đón & 94.021 & $78{,}64\%$\\
Quá thời gian trên xe & 25.300 & $21{,}16\%$\\
Quá sức chứa & 239 & $0{,}20\%$\\
\midrule
Tổng & 119.560 & \\
\bottomrule
\end{tabular}
\end{table}

Tuy nhiên, ``phương án bị loại'' khác với ``hành khách bị từ chối''. Một hành khách chỉ bị từ
chối khi \emph{mọi} phương án cho người đó đều hỏng. Tách riêng 425 hành khách bị từ chối cho
một bức tranh khác hẳn: cả 425 đều có phương án chết vì cửa sổ, và phần vượt cửa sổ \emph{nhỏ
nhất} của mỗi người có trung vị chỉ $44{,}4$ giây, trong khi trung vị trên toàn bộ phương án bị
loại là $309{,}7$ giây. Nói cách khác, đa số phương án hỏng nặng, nhưng phương án tốt nhất của
mỗi hành khách bị từ chối thì chỉ lỡ trong gang tấc.

\subsection{Kết quả phép thử trần}

Để kiểm tra xem cửa sổ hay dung lượng đội xe mới là ràng buộc thật, cửa sổ được nới từ 600 lên
900 giây và hai cấu hình gốc được chạy lại.

\begin{table}[htbp]
\centering
\caption{Phép thử trần: nới cửa sổ đón từ 600 lên 900 giây}
\label{bang:tran}
\small
\begin{tabular}{@{}lrrr@{}}
\toprule
\textbf{Cấu hình} & \textbf{Cửa sổ 600 s} & \textbf{Cửa sổ 900 s} & \textbf{Chênh}\\
\midrule
B1 (không hứa) & 1.735 & 1.938 & $+203$ ($+9{,}40$ pp)\\
C1 (hứa 30 s) & 1.581 & 1.753 & $+172$ ($+7{,}96$ pp)\\
\midrule
Giá của lời hứa & $-7{,}13$ pp & $-8{,}56$ pp & \\
\bottomrule
\end{tabular}
\end{table}

Kết quả bác bỏ giả thuyết ban đầu của chúng tôi. Dự đoán là khi cửa sổ thôi chặn, chi phí của
cam kết sẽ co lại; thực tế nó \emph{nở ra}, từ $7{,}13$ lên $8{,}56$ điểm phần trăm. Cơ chế
tương đối rõ: cửa sổ rộng hơn cho bộ tối ưu nhiều lựa chọn hơn, nên cũng nhiều cơ hội dời lời
hứa hơn, và trần vì thế cắn đau hơn.

Phát biểu tổng quát rút ra được, và nó khớp với con số $+11{,}1$ điểm phần trăm mà Fielbaum và
Alonso-Mora đo được trong một hệ 2.000 xe --- tự do hơn hẳn hệ 8 xe của khóa luận:

\begin{khungy}
Giá của một lời hứa không phải hằng số. Nó tăng theo mức tự do của nhà điều hành.
\end{khungy}

\subsection{Mặt cong hai chiều}
\label{muc:matcong}

Phát biểu trên rút ra từ đúng hai điểm --- cấu hình không cam kết và cấu hình trần 30 giây ---
nên nó là một quan sát chứ chưa phải một quy luật. Để kiểm tra, toàn bộ thang ngân sách được
chạy lại ở cửa sổ 900 giây, cho một mặt cong hai chiều: giá của cam kết như một hàm của cả trần
$\beta$ lẫn mức tự do.

\begin{table}[htbp]
\centering
\caption{Mặt cong giá của cam kết trên hai mức tự do}
\label{bang:matcong}
\small
\begin{tabular}{@{}lrrrrr@{}}
\toprule
\multirow{2}{*}{\textbf{Trần $\beta$}} & \multicolumn{2}{c}{\textbf{Cửa sổ 600 s}} & \multicolumn{2}{c}{\textbf{Cửa sổ 900 s}} & \multirow{2}{*}{\textbf{Chênh}}\\
\cmidrule(lr){2-3}\cmidrule(lr){4-5}
 & phục vụ & vs B1 & phục vụ & vs B1 & \\
\midrule
không trần & 1.655 & $-3{,}70$ & 1.837 & $-4{,}68$ & $-0{,}97$\\
300 s & 1.641 & $-4{,}35$ & 1.835 & $-4{,}77$ & $-0{,}42$\\
120 s & 1.608 & $-5{,}88$ & 1.794 & $-6{,}67$ & $-0{,}79$\\
60 s & 1.579 & $-7{,}22$ & 1.761 & $-8{,}19$ & $-0{,}97$\\
30 s & 1.581 & $-7{,}13$ & 1.753 & $-8{,}56$ & $\mathbf{-1{,}44}$\\
\midrule
B1 (không hứa gì) & 1.735 & --- & 1.938 & --- & \\
\bottomrule
\end{tabular}
\end{table}

Cột cuối âm ở \emph{cả năm} mức trần: nới tự do làm lời hứa đắt hơn ở mọi độ chặt, không chỉ ở
một điểm. Bốn trong năm giá trị nằm rõ trên sàn nhiễu $\pm 0{,}3$ điểm phần trăm; riêng mức
300 giây với $-0{,}42$ pp (9 hành khách) chỉ vừa đủ trên sàn và không nên diễn giải riêng lẻ.

Đáng chú ý hơn là hình dạng. Trên chính thang ngân sách, từ 300 xuống 30 giây, độ chênh tăng
đơn điệu: $-0{,}42 \to -0{,}79 \to -0{,}97 \to -1{,}44$ điểm phần trăm. Trần càng chặt thì việc
cho nhà điều hành thêm tự do càng đắt.

\begin{khungy}
Tự do của nhà điều hành và độ chặt của lời hứa \emph{tương tác} với nhau chứ không cộng độc
lập. Một trần 30 giây tốn thêm $1{,}44$ điểm phần trăm khi cửa sổ nới từ 600 lên 900 giây,
trong khi một trần 300 giây chỉ tốn thêm $0{,}42$ điểm.
\end{khungy}

Cơ chế giải thích khá trực tiếp và khớp với hình dạng đường biên ở mục \ref{muc:duongbien}: cửa sổ rộng hơn cho bộ tối ưu nhiều
phương án hơn, mà mỗi phương án đều là một cơ hội dời lời hứa. Với trần lỏng, phần lớn các cơ
hội đó vẫn lọt; với trần chặt, chúng bị cắt, nên phần tự do tăng thêm bị lãng phí nhiều hơn.

Kết quả này có giá trị thực hành trực tiếp: nó nói rằng không thể chọn trần $\beta$ độc lập với
mức tự do của hệ thống. Một trần hợp lý cho đội xe lớn, mạng thoáng sẽ đắt hơn đáng kể so với
cùng trần đó trên đội xe nhỏ, mạng chật.

Phép đo chất lượng dịch vụ đóng đinh thêm một điểm nữa. Nới cửa sổ làm chờ trung vị của cấu
hình C1 tăng từ 407 lên 636 giây, tức bộ tối ưu tiêu gần hết phần nới thêm vào việc bắt khách
chờ. Quy ra, mỗi hành khách cứu thêm được tốn 2.421 giây chờ cộng thêm trên toàn hệ, trong đó
chỉ 579 giây là của chính người đó --- phần còn lại áp lên người khác. Với hệ số định giá ở mục
2.5, đây là một món hời tồi, và hướng nới cửa sổ đồng loạt bị loại bỏ.

\subsection{Một bất biến miền chặn tắc đường}
\label{muc:batbien}

Thí nghiệm trôi dạt làm lộ ra một giới hạn của chính hệ thống. Khi tắc đường làm xe đang trên
đường đón đến muộn hơn cửa sổ đã chấp nhận lúc nhận khách, hệ thống báo lỗi và dừng. Nguồn gốc
là một bất biến trong tầng miền: khi hành khách lên xe, nếu thời điểm đón nằm ngoài khoảng
$[\text{sớm nhất}, \text{muộn nhất}]$ thì thao tác thất bại. Đây là bất biến \emph{miền} chứ
không phải chính sách --- nó bắn bất kể cơ sở tính nào đang dùng, bất kể nguyên nhân gì.

Đáng chú ý là nới cửa sổ \emph{không} cứu được: thử lại toàn bộ trên bộ dữ liệu cửa sổ 900 giây
vẫn hỏng 20/20 ô với đúng lỗi đó, vì như vừa nêu, bộ tối ưu tiêu hết phần nới. Phát biểu tổng
quát:

\begin{khungy}
Một cửa sổ đón cố định, được chấp nhận trước khi cao điểm hình thành, thì sau đó không giữ nổi
--- bất kể điều phối làm gì. Và nới cửa sổ không tạo ra khoảng đệm, vì bộ tối ưu luôn lập kế
hoạch sát mép bất kỳ cửa sổ nào được cho.
\end{khungy}

Hệ quả cho mô hình: hệ thống hiện chưa biểu diễn được trạng thái mà riêng giao thông đã phá vỡ
một cửa sổ đón đã cam kết. Đây là ranh giới của chính hệ thống, do thực nghiệm lộ ra, và được
ghi nhận như một hạn chế ở mục 5.2.

\section{Đối chiếu với Fielbaum và Alonso-Mora (2020)}
\label{muc:doichieu}

\subsection{Cài đặt thước đo của họ trên dữ liệu của mình}

Vì công trình \cite{fielbaum2020} là hàng xóm gần nhất, khóa luận cài đặt lại chính thước đo
của họ --- công thức (2.5) --- trên biên bản của mình, để so sánh sòng phẳng trên cùng dữ liệu.
Việc này khả thi vì cả hai đại lượng cần thiết đều rút được từ chuỗi lời hứa đã công bố.

\begin{table}[htbp]
\centering
\caption{Thước đo hai đầu của văn liệu so với thước đo chuỗi của khóa luận}
\label{bang:doichieu}
\small
\begin{tabular}{@{}lrr@{}}
\toprule
& \textbf{B1} & \textbf{C1}\\
\midrule
$P(\Delta D > 0)$ & $0{,}2329$ & $0{,}0076$\\
$E(\Delta D \mid \Delta D > 0)$ & $160{,}4$ s & $10{,}5$ s\\
Thước hai đầu (giờ trả) & 64.814 s & 126 s\\
Thước chuỗi của khóa luận (giờ trả) & 71.529 s & 6.822 s\\
\textbf{Tỉ số chuỗi / hai đầu} & $\mathbf{1{,}1\times}$ & $\mathbf{54{,}1\times}$\\
Khách bị chấm 0 dù thực tế bị dời & 663 ($38{,}2\%$) & 774 ($\mathbf{49{,}0\%}$)\\
\bottomrule
\end{tabular}
\end{table}

Ba điều rút ra, và thứ tự của chúng quan trọng.

\textbf{Trên B1, hai thước gần như trùng nhau.} Tỉ số $1{,}1$ lần. Đây là một phép \emph{kiểm
định chéo} cho thước đo của khóa luận: nó không bịa ra chuyển động. Trong chế độ không ràng buộc,
lời hứa trôi đơn điệu ra xa nên hiệu hai đầu xấp xỉ độ dài đường đi.

\textbf{Trên C1, hai thước lệch 54 lần.} Vì chế độ cam kết ép bộ tối ưu kéo lời hứa quay về gần
giá trị đã công bố, mà thước hai đầu chỉ nhìn hai điểm nên không thấy đường đi. Nói cách khác,
thước đo của văn liệu \emph{mù đúng ở chế độ mà khóa luận nghiên cứu}.

\textbf{Thước đếm của họ không phân biệt được hai cấu hình.} Số lần dự báo tăng trung bình là
$14{,}99$ ở B1 và $13{,}45$ ở C1, trong khi hai cấu hình cách nhau $7{,}13$ điểm phần trăm dịch
vụ. Tuy nhiên phải kèm cảnh báo: hệ thống của khóa luận công bố lại lời hứa khoảng 30 lần mỗi
hành khách, còn hệ của họ cập nhật mỗi phút nên chỉ khoảng 5 lần. Chỉ được so B1 với C1 trong
nội bộ dữ liệu của mình, không được so trực tiếp với con số in trong bài của họ.

\subsection{Hình dạng của kết quả}

Cặp $1{,}1\times$ và $54{,}1\times$ có một hình dạng đáng chú ý: thước đo của khóa luận
\emph{quy về} thước đo của văn liệu trong chính chế độ mà văn liệu nghiên cứu, và chỉ tách ra ở
chế độ họ chưa có. Đó là quan hệ giữa một trường hợp tổng quát và một trường hợp riêng, chứ
không phải hai thước đo mâu thuẫn. Nếu chúng lệch nhau ở mọi chế độ thì mới đáng ngờ.

\subsection{Đối chiếu ở mức cơ chế}

Ngoài thước đo, khóa luận cũng cài đặt lại \emph{cơ chế} của họ. Kịch bản ``kiểm soát hoàn
toàn'' trong bài của họ --- đánh dấu hành khách là không thể bị từ chối và siết thời gian chờ
tối đa xuống bằng đúng giờ đã báo --- chính là cơ chế ratchet ở mục \ref{muc:ratchet}, và về bản chất là cơ
sở kết quả với trần bằng không.

\begin{table}[htbp]
\centering
\caption{So sánh tương đối giá của việc kiểm soát}
\label{bang:sosanh-tuongdoi}
\small
\begin{tabular}{@{}lrrrr@{}}
\toprule
& \textbf{Gốc} & \textbf{Siết chặt} & \textbf{Tuyệt đối} & \textbf{Tương đối}\\
\midrule
Fielbaum \& Alonso-Mora (2020) & $34{,}5\%$ & $45{,}6\%$ & $+11{,}1$ pp & $1{,}32\times$\\
Khóa luận này (B1 $\to$ C1) & $19{,}68\%$ & $26{,}81\%$ & $+7{,}13$ pp & $1{,}36\times$\\
\bottomrule
\end{tabular}
\end{table}

Đọc bảng này cần trung thực. Tính theo điểm phần trăm thì cơ chế của khóa luận rẻ hơn, nhưng
tính theo \emph{tỉ lệ} thì nhỉnh hơn một chút về phía xấu. Đây là một \emph{tái lập độc lập}
chứ không phải một chiến thắng: trên một thực thể nhỏ hơn 185 lần, cùng cơ chế cho cùng kết luận
định tính với hệ số gần như trùng. Với tiêu chí minh chứng định lượng, tái lập trên thực thể độc
lập là điểm cộng thật; nó không phải là bằng chứng vượt trội.

\subsection{Phần trùng và phần không trùng}

Để tránh mọi hiểu lầm, bảng \ref{bang:trung} liệt kê rõ ranh giới.

\begin{table}[htbp]
\centering
\caption{Ranh giới với công trình gần nhất}
\label{bang:trung}
\small
\begin{tabular}{@{}p{7.2cm}cc@{}}
\toprule
\textbf{Nội dung} & \textbf{2020} & \textbf{Khóa luận}\\
\midrule
Khái niệm lời hứa công bố bị xê dịch & có & có\\
Đo độ lớn dịch chuyển trên từng hành khách & có & có\\
Đếm số lần dự báo bị đổi & có & có\\
Phân loại nguyên nhân (định tính, bằng lời) & có & có\\
Đường cong đánh đổi ổn định--phục vụ & có (3 điểm) & có (9 điểm)\\
Cơ chế ratchet trên giờ đón & có & có\\
\midrule
Trần số cưỡng chế được trên xê dịch cộng dồn & \textbf{không} & \textbf{có}\\
Dạng cộng dồn theo chuỗi (độ lớn $\times$ chuỗi) & \textbf{không} & \textbf{có}\\
Phân rã nhân quả định lượng theo từng quyết định & \textbf{không} & \textbf{có}\\
Phân biệt hai cơ sở tính và ba định lý & \textbf{không} & \textbf{có}\\
\bottomrule
\end{tabular}
\end{table}

Sáu dòng đầu là phần trùng, và khóa luận coi đó là \emph{nền} chứ không phải đóng góp. Bốn dòng
sau là phần khác biệt. Đáng chú ý là chính các tác giả kết thúc mục 6 của bài mình bằng câu
``điều này càng cho thấy cần những ý tưởng tinh vi hơn để kiểm soát bất ổn'' --- tức là họ gọi
tên chính chỗ trống mà khóa luận này lấp vào.

\section{Phúc lợi: điểm hòa vốn}
\label{muc:phucloi}

Mục này quy đổi kết quả kỹ thuật sang ngôn ngữ phúc lợi, với đầy đủ cảnh báo.

Lượng thời gian bất ngờ mà cam kết tránh được là 74.255 giây, tức $20{,}63$ giờ. Số chuyến phải
hy sinh để tránh được nó là 154. Vậy mỗi chuyến bị hy sinh đổi lấy 482 giây --- khoảng tám phút
--- thời gian bất ngờ tránh được cho tập thể.

\begin{table}[htbp]
\centering
\caption{Điểm hòa vốn theo mức định giá thời gian bất ngờ}
\label{bang:hoavon}
\small
\begin{tabular}{@{}lr@{}}
\toprule
\textbf{Định giá thời gian bất ngờ} & \textbf{Giá trị hòa vốn của một chuyến}\\
\midrule
$7{,}88$ EUR/h (giá trị thời gian trên xe, cận dưới) & $1{,}06$ EUR\\
$10{,}80$ EUR/h (cận trên) & $1{,}45$ EUR\\
$20{,}00$ EUR/h (chờ bất ngờ, cận dưới) & $2{,}68$ EUR\\
$25{,}00$ EUR/h (chờ bất ngờ, cận trên) & $3{,}35$ EUR\\
\bottomrule
\end{tabular}
\end{table}

Cách phát biểu trung thực là: với mức định giá thời gian bất ngờ đã công bố trong văn liệu, hệ
thống có cam kết vượt hệ thống không cam kết về phúc lợi khi thặng dư của một chuyến ghép biên
thấp hơn khoảng $2{,}7$--$3{,}4$ EUR. Vì những chuyến bị từ chối là những chuyến khó phục vụ
nhất, tức gần điểm bàng quan nhất, thặng dư biên của chúng theo định nghĩa là nhỏ.

Đây là một \emph{ngưỡng}, không phải một tuyên bố chiến thắng. Ba giả định cần ghi rõ: phép tính
coi mọi chuyến bị mất có cùng giá trị thặng dư, đó là đơn giản hóa; đại lượng đo là tổng dịch
chuyển tích lũy chứ không phải sai số cuối cùng, nên phải gọi đúng tên là \emph{độ xáo trộn của
lời hứa} chứ không phải \emph{độ trễ}; và như mục 2.5 đã nêu, giá trị độ tin cậy đo được trong
bối cảnh này thấp hơn giá trị thời gian, nên lập luận phúc lợi ở đây kém chắc chắn hơn phần kỹ
thuật khá nhiều.

\section{Chứng cứ về tính không giả tạo}
\label{muc:chungcu}

Một phản biện tự nhiên với mọi kết quả dạng ``ràng buộc làm giảm hiệu năng'' là: có thể hiệu
năng giảm không phải vì ràng buộc mà vì tìm kiếm bị cắt ngắn. Khóa luận kiểm tra điều này trực
tiếp trên toàn bộ 40 lần chạy của hai cấu hình gốc.

\begin{table}[htbp]
\centering
\caption{Kiểm tra tìm kiếm không bị cắt ngắn}
\label{bang:khonggiatao}
\small
\begin{tabular}{@{}lr@{}}
\toprule
\textbf{Chỉ báo} & \textbf{Giá trị}\\
\midrule
Tỉ lệ lần quyết định cạn ngân sách tìm kiếm & $0\%$\\
Tỉ lệ lần quyết định chạm trần số ứng viên & $0\%$\\
Số phương án khả thi bị bỏ sót do trần & 0\\
Số nhánh chưa được mở rộng & 0\\
Tỉ lệ loại bỏ do ràng buộc vật lý & $94\%$\\
Tỉ lệ loại bỏ do ràng buộc cam kết & $1{,}1\%$\\
\bottomrule
\end{tabular}
\end{table}

Không gian tìm kiếm chưa bao giờ bị cắt ngắn, và ràng buộc cam kết chỉ chiếm $1{,}1\%$ tổng số
lần loại bỏ --- một ràng buộc \emph{hẹp nhưng đòn bẩy cao}. Nghĩa là chênh lệch tỉ lệ phục vụ
giữa các cấu hình phản ánh đúng tác động của cam kết, không phải sản phẩm phụ của giới hạn tính
toán.

\section{Bộ sinh phương án biết đến lời hứa}
\label{muc:bosinh}

Toàn bộ kết quả từ đầu chương đến đây đều dùng một bộ sinh phương án \emph{không} biết gì về
lời hứa: nó sinh ra các phương án theo tiêu chí vật lý và lịch trình, rồi cổng cam kết mới xóa
đi những phương án vi phạm. Câu hỏi tự nhiên tiếp theo là liệu có thu được gì nếu bộ sinh biết
trước về ngân sách, thay vì để cổng dọn dẹp phía sau. Mục này trả lời bằng phép đo, và câu trả
lời phần lớn là không --- nhưng lý do \emph{tại sao} không lại là một kết quả có giá trị riêng.

\subsection{Ba mức can thiệp và điều kiện để chúng có ý nghĩa}

Chúng tôi tách ý tưởng ``sinh phương án có ý thức về cam kết'' thành ba mức, xếp theo mức độ
xâm lấn tăng dần, vì mỗi mức đòi hỏi một điều kiện khác nhau của dữ liệu để có tác dụng.

\textbf{Mức A --- xếp hạng vị trí chèn.} Bộ sinh vẫn sinh đủ, nhưng ưu tiên những vị trí chèn
làm tiêu ít ngân sách nhất. Mức này chỉ có ý nghĩa nếu mỗi quyết định thực sự có nhiều hơn một
phương án để xếp hạng.

\textbf{Mức B --- cắt tỉa sớm.} Bộ sinh tính một cận dưới cho chi phí ngân sách của cả một
nhánh và bỏ nhánh đó ngay nếu cận dưới đã vượt trần. Mức này chỉ có ý nghĩa nếu hiện có phương
án bị bỏ sót vì hết thời gian hoặc hết chỗ, hoặc nếu tốc độ đang là vấn đề.

\textbf{Mức C --- đặt trước ngân sách.} Không động vào bộ sinh; thay vào đó chèn một tầng
\emph{cảnh báo} vào hàm mục tiêu từ điển của bài toán đóng gói tập hợp, để một phương án còn
trong trần cứng nhưng đã vượt ngưỡng cảnh báo thấp hơn sẽ bị hạ bậc. Ý đồ là làm hệ thống tiêu
ngân sách một cách dè dặt, để dành chỗ cho yêu cầu đến sau.

\subsection{Tập lựa chọn mỏng đến mức nào}

Điều kiện của mức A và mức B đều là mệnh đề về dữ liệu, nên đo được trước khi viết một dòng mã
nào. Hệ thống đã ghi sẵn chẩn đoán cần thiết cho mỗi xe ở mỗi vòng quyết định, trong nhánh
\texttt{executionEvidence/generation/vehicleLosses}. Bảng \ref{bang:mong} tổng hợp trên toàn bộ
panel 20 ô đã niêm phong, cả hai nhánh, tổng cộng $217.736$ quyết định cấp xe.

\begin{table}[htbp]
\centering
\caption{Độ mỏng của tập lựa chọn trên toàn panel niêm phong}
\label{bang:mong}
\small
\begin{tabular}{@{}lrr@{}}
\toprule
\textbf{Đại lượng} & \textbf{B1 (không hứa)} & \textbf{C1 (trần 30 s)}\\
\midrule
Số quyết định cấp xe & 113.064 & 104.672\\
Ứng viên khả thi trước trần / xe & 1,044 & 1,047\\
Ứng viên được giữ lại / xe & 1,044 & 1,047\\
Bị loại vì vật lý hoặc lịch trình / xe & 1,057 & 0,956\\
\midrule
Số lần chạm trần số ứng viên & 0 & 0\\
Số lần cạn ngân sách tính toán & 0 & 0\\
Ứng viên khả thi bị bỏ do trần & 0 & 0\\
Nhánh chưa khai triển bị bỏ & 0 & 0\\
\midrule
Quyết định chỉ có \emph{một} ứng viên & $96{,}04\%$ & $95{,}87\%$\\
Quyết định có từ hai ứng viên trở lên & $3{,}96\%$ & $4{,}13\%$\\
\bottomrule
\end{tabular}
\end{table}

Hai dòng quan trọng nhất là hai dòng bằng nhau: số ứng viên trước trần và số ứng viên sau trần
giống hệt nhau tới ba chữ số thập phân. Bộ giữ ứng viên có cài đặt một trần số lượng, nhưng
trần ấy \emph{chưa bao giờ} chạm trong hai trăm nghìn quyết định.

Hệ quả rất trực tiếp. Mức A là phép toán rỗng: với một ứng viên thì không có gì để xếp hạng, và
việc xếp hạng chỉ có thể tác động lên khoảng bốn phần trăm số quyết định, nơi mà chín phần mười
trường hợp tập chỉ có đúng hai phần tử. Mức B cũng là phép toán rỗng: không có ứng viên khả thi
nào từng bị bỏ, không có nhánh nào chưa được khai triển, nên cắt tỉa sớm không tiết kiệm được
gì; và không có lần nào cạn ngân sách tính toán, nên cũng không có vấn đề tốc độ để giải quyết.

Điều này cũng trả lời một câu hỏi bỏ ngỏ từ chương 3: các kết quả của khóa luận có phải là hệ
quả của việc tìm kiếm bị cắt ngắn hay không. Câu trả lời là không --- tìm kiếm chưa bao giờ bị
cắt. Ràng buộc siết chặt nằm ở tầng vật lý, đúng như phân tích nguyên nhân bị loại đã chỉ ra.

\subsection{Bộ sinh biết sắp xếp lại: một bất biến thứ hai}

Mục 4 của danh sách hướng mở là một bộ sinh có thể \emph{sắp xếp lại} thứ tự các điểm dừng của
hành khách cũ, chứ không chỉ chèn thêm vào giữa. Hóa ra hệ thống đã có sẵn cơ chế ấy: một bộ
dựng mầm sửa chữa lấy một hành khách đang chờ, gỡ cả cặp đón--trả của người đó ra khỏi đuôi khả
biến, rồi chèn lại cặp ấy vào mọi vị trí còn giữ được thứ tự đón trước trả. Một nhánh thực
nghiệm với hạn mức bốn hành khách mỗi xe đã được chạy trên cùng panel, cùng bộ chạy, cùng hạt
giống.

Trước hết, cơ chế này \emph{thực sự} sắp xếp lại, và đó không phải điều hiển nhiên: bộ sinh chèn
thuần chưa một lần đảo thứ tự trong toàn bộ dữ liệu của khóa luận. Bảng \ref{bang:sapxep} cho
thấy sự khác biệt.

\begin{table}[htbp]
\centering
\caption{Hành vi của bộ sinh chèn thuần và bộ sinh biết sắp xếp lại}
\label{bang:sapxep}
\small
\begin{tabular}{@{}lrrrrr@{}}
\toprule
\textbf{Nhánh} & \textbf{Ô} & \textbf{Cập nhật} & \textbf{Sắp lại} & \textbf{Đẩy chặng đầu} & \textbf{Đóng băng sớm}\\
\midrule
B1 --- chèn thuần & 20 & 5.328 & 0 & 57 \ ($1{,}07\%$) & 0\\
C1 --- chèn thuần, trần 30 s & 20 & 4.864 & 0 & 17 \ ($0{,}35\%$) & 0\\
b4 --- biết sắp xếp lại & 20 & 1.978 & 245 & 150 \ ($7{,}58\%$) & 0\\
\bottomrule
\end{tabular}
\end{table}

Chiều \texttt{incumbentOrderInversionCount} của vector ngân sách, vốn luôn bằng không trong mọi
thực nghiệm trước, nay bị tính phí 116 lần với tổng 120 đơn vị. Cơ chế hoạt động đúng như thiết
kế.

Nhưng nhánh này \textbf{hỏng trên 18 trên 20 ô}, với cùng một mã lỗi:
\texttt{LOCKED\_LEG\_MISMATCH} tại chặng đang chạy của một xe. Nguyên nhân là một bất biến thứ
hai của mô hình, khác với bất biến cửa sổ đón ở mục \ref{muc:batbien}, và đáng ghi lại vì nó
phản ánh một bất đồng khái niệm chứ không phải một lỗi cài đặt.

Phía mô hình miền, một điểm dừng chuyển từ đuôi khả biến sang tiền tố đóng băng \emph{khi xe
tới nơi}. Không có chuyển trạng thái nào cho tình huống ``xe đã xuất phát hướng tới''. Phía bộ
mô phỏng, chặng đang chạy bị khóa ngay khi xe lăn bánh, và mọi kế hoạch thay thế có chặng đầu
khác đi đều bị từ chối. Giữa lúc xuất phát và lúc tới, điểm dừng đích nằm ở vị trí đầu của đuôi
khả biến, và bộ sinh biết sắp xếp lại coi nó là sửa được.

Cột cuối của bảng \ref{bang:sapxep} xác nhận điều này ở dạng trực tiếp nhất: trên $12.170$ lần
cập nhật kế hoạch của cả ba nhánh, hệ thống \emph{chưa một lần} đóng băng một điểm dừng trước
khi xe tới, mặc dù kiểu dữ liệu lộ trình cho phép trạng thái đó. Bộ sinh chèn thuần không bao
giờ làm lộ bất đồng này vì nó đẩy chặng đầu rất hiếm --- khoảng một phần trăm số lần cập nhật.
Bộ sinh biết sắp xếp lại đẩy chặng đầu bảy phần trăm, và với khoảng một trăm lần cập nhật mỗi
ô thì việc bắt gặp một chiếc xe đang giữa đường gần như chắc chắn xảy ra.

Chỗ đáng suy nghĩ là bất biến này nói cùng một điều với luận điểm trung tâm của khóa luận, chỉ
trên một trục khác. Phép phân rã ba mốc ở chương 3 vạch ranh giới giữa phần nhà điều hành còn
đổi được và phần thế giới vật lý đã cố định, và nó vạch ranh giới ấy \emph{theo thời gian}: mốc
ngoại sinh tính chính xác kế hoạch cũ sẽ ra sao nếu không ai can thiệp. Tiền tố đóng băng vạch
đúng ranh giới ấy \emph{theo cấu trúc lộ trình}: điểm dừng nào đã hết sửa được. Và phát hiện ở
đây là hệ thống đang vạch ranh giới thứ hai sai chỗ --- ở lúc xe tới thay vì lúc xe xuất phát.
Sai lệch ấy vô hại với bộ sinh chèn thuần và chỉ lộ ra khi có một cơ chế thật sự muốn sắp xếp
lại. Nói cách khác, hai bất biến chặn hai hướng cải tiến ở mục này và mục \ref{muc:batbien} đều
là cùng một câu hỏi ``cái gì đã ngoài tầm kiểm soát'', đặt một lần theo thời gian và một lần
theo thứ tự.

Trên ô duy nhất chạy trọn vẹn, kết quả cũng đáng ghi: nhánh biết sắp xếp lại phục vụ 99 hành
khách, đúng bằng B1, trong khi C1 phục vụ 92. Cần nói rõ rằng nhánh này chạy với toàn bộ trần
cam kết bị tháo bỏ --- hàm chọn chính sách đưa nó về hồ sơ mục tiêu chi phí thuần và dựng lại
mọi chính sách với trần rỗng --- nên nó so được với B1 chứ không so được với C1. Đọc như vậy thì
kết luận là: thêm khả năng sắp xếp lại vào một bộ điều phối vốn đã tự do hoàn toàn thu được đúng
không hành khách nào. Cỡ mẫu một ô nên đây là chỉ dấu chứ chưa phải kết luận, nhưng nó phù hợp
với phát hiện trung tâm rằng ràng buộc siết chặt là cửa sổ đón chứ không phải thứ tự điểm dừng.

Một câu hỏi tự nhiên là liệu việc đóng băng chặng đang chạy có làm cơ chế sắp xếp lại mất hết ý
nghĩa hay không. Phân loại $245$ lần sắp xếp lại quan sát được cho thấy không: $102$ lần
($41{,}6\%$) diễn ra hoàn toàn phía sau chặng đang chạy và vẫn hợp lệ dưới một phép đóng băng
như vậy; $134$ lần ($54{,}7\%$) là đúng những nước đi đang làm sập chương trình; $9$ lần còn lại
rơi vào trường hợp chặng đầu vừa được phục vụ xong. Nghĩa là bước sửa này lấy đi khoảng một nửa
láng giềng nhưng giữ lại một nửa --- đủ để câu hỏi còn đáng hỏi.

Cuối cùng, tổ hợp mà mục 4 thực sự cần --- sắp xếp lại \emph{cộng với} lời hứa có biên --- hiện
không diễn đạt được: bộ đọc cấu hình từ chối trường sửa chữa khi chính sách không phải nhánh
b4. Muốn trả lời câu hỏi ấy phải sửa mã nguồn ở ba chỗ theo đúng thứ tự: thêm chuyển trạng thái
đóng băng chặng đang bay ở tầng miền, cho chính sách vector cứng nhận hạn mức sửa chữa, và nới
điều kiện kiểm tra cấu hình. Chúng tôi chuyển việc này sang chương hướng phát triển, kèm ghi chú
rằng bằng chứng ở đoạn trên gợi ý phần thưởng nhiều khả năng bằng không.

\subsection{Đặt trước ngân sách: chạy được nhưng không đổi được gì}

Mức C là hạng mục duy nhất chạy được ngay trên hệ thống hiện có, vì cơ chế cần thiết đã được
cài đặt: một hồ sơ mục tiêu lai mềm--cứng chèn thêm một tầng \emph{cảnh báo} vào giữa, cho ra
thứ tự từ điển

\begin{center}
\texttt{worst-hard-utilization} $\;>\;$ \texttt{warning-excess} $\;>\;$
\texttt{revision} $\;>\;$ \texttt{cost}
\end{center}

Ngưỡng cảnh báo bắt buộc phải thấp hơn trần cứng, nên phép ghép khả dụng là ngưỡng 60 giây với
trần 120 giây. Nhánh đối chứng --- cùng trần 120 giây nhưng không có tầng cảnh báo --- đã có
trong thang ngân sách ở mục \ref{muc:duongbien}, nên chỉ cần chạy thêm 20 ô nhánh điều trị trên
cùng bộ chạy, cùng fixture, cùng hạt giống.

Kết quả là hai nhánh \textbf{trùng khớp tới từng quyết định} trên cả hai mươi ô. Cả hai phục vụ
$1.608$ trên $2.160$ hành khách, tức giá lời hứa $-5{,}88$ điểm phần trăm ở cả hai nhánh, không
một ô nào lệch dù một hành khách. Nhưng điểm đáng nói không nằm ở con số phục vụ: trùng khớp cả
ở chuỗi kế hoạch được gán cho từng xe, ở tập lời hứa công bố trong từng khung quyết định, ở chuỗi
mã lý do, và ở cả chín thước đo bất ổn của mục \ref{muc:dovung} --- $47.118$ lần công bố, tổng
xáo trộn do quyết định $5.063{,}3$ giây, hành khách bị dịch nặng nhất $119{,}0$ giây, tất cả
giống nhau tới từng chữ số. Bám ngữ nghĩa niêm phong của hai nhánh \emph{có} khác nhau, nhưng
khác vì nhánh lai ghi thêm chứng cứ cảnh báo và phát các lần công bố theo thứ tự khác bên trong
cùng một khung --- không vì một quyết định nào đổi.

Câu hỏi thú vị là \emph{vì sao}. Một tầng trong hàm mục tiêu từ điển có thể vô hiệu theo hai
cách rất khác nhau: giá trị của nó luôn bằng không, tức ngưỡng chưa bao giờ bị chạm; hoặc giá
trị có biến thiên nhưng một tầng cao hơn đã phân định xong mọi so sánh. Bảng \ref{bang:tang},
đọc từ chứng cứ solver, cho thấy đây là trường hợp thứ hai.

\begin{table}[htbp]
\centering
\caption{Các tầng mục tiêu của bài toán đóng gói tập hợp ở nhánh mức C}
\label{bang:tang}
\small
\begin{tabular}{@{}lrrr@{}}
\toprule
\textbf{Tầng mục tiêu} & \textbf{Số lần} & \textbf{Khác 0} & \textbf{Tỉ lệ}\\
\midrule
\texttt{accepted-request-count} & 4.264 & 519 & $12{,}2\%$\\
\texttt{operational-cost} & 4.264 & 4.240 & $99{,}4\%$\\
\texttt{worst-hard-utilization-ppm} & 4.240 & 4.240 & $\mathbf{100{,}00\%}$\\
\texttt{warning-excess:drop\_eta\_total\_ms} & 4.240 & 810 & $19{,}1\%$\\
chín chiều cảnh báo còn lại & 4.240 & 0 & $0\%$\\
\texttt{revision:drop\_eta\_total\_ms} & 4.240 & 21 & $0{,}5\%$\\
\bottomrule
\end{tabular}
\end{table}

Ngưỡng 60 giây \emph{có} bị chạm, trong $19{,}1\%$ số quyết định. Nhưng tầng ngay phía trên nó
khác không trong \emph{toàn bộ} số quyết định. Trong một thứ tự từ điển, một tầng chỉ được lên
tiếng khi mọi tầng trên nó đã hòa; tầng trên ở đây thì hầu như không bao giờ hòa, vì nó đo mức
tiêu ngân sách tệ nhất theo phần triệu --- một thang \emph{liên tục} của đúng cái đại lượng mà
ngưỡng cảnh báo chỉ chia thành hai bậc.

Đây là kết quả mà chúng tôi cho là có giá trị hơn một phép thử thất bại thông thường: ý tưởng
đặt trước ngân sách không sai, nó \emph{đã} nằm trong hệ thống ở dạng mạnh hơn. Cực tiểu hóa mức
tiêu tệ nhất chính là ``tiêu dè dặt'', và nó thực hiện việc đó mịn hơn bất kỳ ngưỡng nào. Một
tầng ngưỡng đặt bên dưới một tầng liên tục cùng ý nghĩa là dư thừa về mặt cấu trúc, không phải
vì dữ liệu không thuận. Muốn mức C có tác dụng thật thì phải hoặc đưa tầng cảnh báo lên
\emph{trên} tầng liên tục, hoặc bỏ tầng liên tục khỏi hồ sơ lai --- cả hai đều làm hệ thống tệ
hơn theo chính thước đo mà khóa luận dùng, nên chúng tôi không theo.

\section{Các hướng đã thử và lý do từ bỏ}
\label{muc:dathu}

Phần này ghi lại những hướng đã được cân nhắc nghiêm túc và bị loại bỏ bằng dữ liệu. Chúng tôi
giữ lại vì tin rằng chúng có giá trị không kém những hướng đã thành công, và vì việc loại bỏ
sớm bằng phép đo rẻ tiền là một phần của phương pháp.

\textbf{Đệm ETA đồng loạt.} Ý tưởng là báo cho khách một giờ muộn hơn thực tế để có chỗ hấp thụ
biến động. Phép đo cho thấy cần tổng cộng 265.000 giây đệm để tránh 22.485 giây xáo trộn --- tỉ
lệ 12 trên 1. Loại bỏ.

\textbf{Hứa theo hạn chót thay vì theo điểm.} Thay vì hứa ``đón lúc 9:10'', hứa ``đón trước
9:15''. Phân tích cho thấy đây chính xác là cửa sổ $\Omega$ và $\Delta$ đã có sẵn, nên không
phải cơ chế mới.

\textbf{Sinh phương án có ý thức về cam kết.} Ý tưởng là để bộ sinh tránh ngay từ đầu những
phương án sẽ vượt trần. Phép đo phân bố mức vượt trần cho thấy phân vị 10 đã là 31 giây trên
trần 30 giây, tức phần lớn phương án vượt rất xa; nới trần lên 60 giây chỉ thu hồi được
$9{,}9\%$. Mục \ref{muc:bosinh} tách ý tưởng này thành ba mức và bác bỏ hai mức đầu bằng phép
đo trên toàn panel.

\textbf{Sắp xếp lại thứ tự hành khách cũ.} Đo được \emph{không một lần} đảo thứ tự nào trong
toàn bộ dữ liệu của bộ sinh chèn thuần, và một phép thử mở rộng láng giềng gấp $8{,}3$ lần cho
$98{,}63\%$ trường hợp vô nghiệm và không thu hồi được hành khách nào. Mục \ref{muc:bosinh}
kiểm tra tiếp bằng chính cơ chế sắp xếp lại đã có sẵn trong hệ thống, và tìm ra một bất biến
thứ hai chặn hướng này lại.

\textbf{Cổng mềm thay cho cổng cứng.} Ý tưởng là thay vì xóa phương án vượt trần, chỉ phạt điểm.
Kiểm tra mã nguồn cho thấy cơ chế hiện có \emph{không} gỡ cổng cứng: phương án chỉ được chấm
điểm cảnh báo \emph{sau khi} đã qua cổng cứng, nên cơ chế thu hồi mà chúng tôi kỳ vọng không tồn
tại. Ngoài ra, một cổng mềm thật sự sẽ biến lời hứa từ cưỡng chế được thành thống kê, tức đánh
đổi mất Định lý 3.1 --- tài sản chính của khóa luận.

\textbf{Cửa sổ đón thương lượng có chọn lọc.} Đây là hướng còn sống sót lâu nhất và được bàn kỹ
nhất. Ý tưởng là khi một hành khách sắp bị từ chối vì lỡ cửa sổ, thay vì từ chối thì thương
lượng một lần rồi khóa lời hứa. Phép thử trần ở mục \ref{muc:cuaso} cho thấy trần dịch vụ tồn tại
($+172$ hành khách) nhưng cách nới đồng loạt lấy phần đó bằng cách đẩy mọi người chờ thêm 229
giây. Cách nới có chọn lọc theo định nghĩa không làm vậy, nên hướng này \emph{chưa} bị bác bỏ;
tuy nhiên nó đòi hỏi sửa mã nguồn và mở rộng phát biểu vấn đề sang kiểm soát tiếp nhận, nên
được chuyển sang chương hướng phát triển.

\section{Hạn chế của thực nghiệm}
\label{muc:hanche}

Cuối chương, chúng tôi liệt kê các hạn chế đã biết.

\textbf{Quy mô nhỏ.} 108 yêu cầu và 8 xe mỗi ô là nhỏ hơn nhiều so với các thực nghiệm quy mô
Manhattan. Các con số tuyệt đối không so trực tiếp được với văn liệu; chỉ so sánh nội bộ giữa
các cấu hình mới có ý nghĩa.

\textbf{Trôi dạt ở panel chính là cận dưới.} Như mục \ref{muc:troidat} đã nêu, panel gốc có đồ thị thời gian
di chuyển tĩnh. Kết luận ``hai cơ sở gần như không khác nhau về số hành khách'' là tính chất của
panel đó, không phải của bài toán; thí nghiệm bơm trôi dạt đã bù một phần nhưng chưa trọn vẹn.

\textbf{Không có kết quả về tỉ lệ phục vụ dưới trôi dạt thật.} Bất biến miền ở mục \ref{muc:batbien} chặn phần
này.

\textbf{Hàng rào dài hơn tầm đặt chỗ chưa kiểm chứng được.} Vì tầm đặt chỗ đúng 600 giây, mọi
hàng rào từ mức đó trở lên đều tương đương khóa vĩnh viễn.

\textbf{Lập luận phúc lợi kém chắc chắn hơn phần kỹ thuật.} Nó phụ thuộc vào hệ số định giá từ
một bối cảnh địa lý khác và vào giả định về thặng dư biên của chuyến bị mất.

\textbf{Chưa cho phép tái gán hành khách sang xe khác.} Mọi cấu hình đều cấm điều này, nên con
số quy trách $89{,}06\%$ chưa tính phần do đổi xe.
